\section{Objetivos}

\textbf{Objetivo General}

Desarrollar un marco de aprendizaje estocástico y generalizable basado en modelos fundacionales para el análisis de señales electroencefalográficas (EEG), que permita aprovechar datos no etiquetados, manejar la variabilidad entre bases de datos y proporcionar estimaciones de la incertidumbre para apoyar el diagnóstico de trastornos mentales.

\textbf{Objetivos específicos}

\begin{enumerate}
	\item Diseñar un enfoque de aprendizaje auto-supervisado para señales EEG que permita extraer representaciones latentes transferibles a partir de datos no etiquetados, reduciendo la dependencia de conjuntos de datos anotados.
	
	\item Desarrollar una estrategia de generalización de dominio que permita mejorar la robustez del modelo frente a la variabilidad entre bases de datos EEG con diferentes protocolos de adquisición, incluyendo diferencias en tasas de muestreo y número de canales.
	
	\item Integrar un mecanismo de inferencia estocástica en el espacio de representaciones aprendidas que permita estimar la incertidumbre de las predicciones y evaluar su desempeño, con el fin de mejorar la confiabilidad del modelo en escenarios clínicos.
\end{enumerate}